\section{Object, Thesis, and Rules of Reading}

\begin{studybox}{Governing thesis}
The Roman Order of Mass is one covenantal action: the baptized are gathered and purified; God addresses his Church; the Church confesses, offers, and gives thanks through the ordained ministry acting in Christ's person; Christ's one sacrifice becomes sacramentally present; the faithful are ordered to Communion; and the blessing sends the worshipping body into mission. Its many historical strata are not thirty-three independent devotions. Each part receives and transforms what precedes it, and the whole is governed by the Paschal passage of Christ to the Father in the Holy Spirit.
\end{studybox}

This study takes \emph{Ordinary} in its broad sense: the stable spoken and sung texts and the complete ritual framework printed in the 1962 \latin{Missale Romanum}. It therefore exceeds the five choral items often called the Mass Ordinary---Kyrie, Gloria, Credo, Sanctus, and Agnus Dei---and follows the complete \latin{Ordo Missae}, including gestures, speakers, silence, and the stable positions occupied by variable propers. It distinguishes Low, sung, and Solemn Mass when differentiated ministry, audibility, procession, or incense materially changes how a unit appears.

The number thirty-three is an editorial map, not a rubric or article of faith. It has a documented modern catechetical antecedent: Msgr John T. McMahon's \emph{Pray the Mass} (13th revised edition, 1959) proposes thirty-three ``steps,'' explicitly one for each traditionally reckoned year of Christ's earthly life. That scheme fruitfully invites the worshipper to read the Mass christologically, but it is uneven as a complete ritual anatomy: changing the Missal is one step, while several distinct postcommunion acts are gathered together. This study therefore acknowledges McMahon without reproducing him. The Missal supplies the sequence; the present inventory redraws thirty-three boundaries at comprehensive textual and ritual joints so that every part can be studied and then returned to the whole. The numerical association may remain spiritually suggestive, but it neither establishes the divisions nor proves an original Roman intention. Catalog identifier \texttt{O00} likewise means ``read this stable framework before the variable propers''; the Ordinary is not itself a proper and is not a fictitious Sunday.

\subsection{Five layers that must not be confused}

\begin{longtable}{>{\bfseries\raggedright\arraybackslash}p{0.19\linewidth}>{\raggedright\arraybackslash}p{0.73\linewidth}}
Received order & The Latin text, sequence, minister, gesture, audibility, and conditions attested by the 1962 Missal.\\
Biblical and doctrinal reality & What Scripture and the Church's teaching establish about sacrifice, memorial, real presence, priesthood, Communion, and ecclesial worship.\\
Historical stratum & The period and route by which a text or gesture entered the Roman order, stated with the uncertainty allowed by the evidence.\\
Received spiritual interpretation & Patristic, medieval, or later Catholic readings that genuinely illuminate the rite, even when they are not its first compositional explanation.\\
Editorial synthesis & This study's account of how units relate. Such synthesis is offered for testing and never placed in an authority's mouth.\\
\end{longtable}

The distinction protects two truths at once. The Mass has a divinely instituted sacramental nucleus: Christ gave the Church his Body and Blood under bread and wine and commanded the apostolic ministry to do this in his memorial. Its complete Roman ceremonial form also developed historically under ecclesial authority. Organic growth does not make every medieval prayer apostolic; historical addition does not make the received rite arbitrary or doctrinally mute. Pius XII describes the liturgy as the public worship of the whole Mystical Body and also as a living reality that develops while the Church guards faith and orders worship (\emph{Mediator Dei} 20--23, 47--59).

\subsection{The act beneath the inventory}

The classic distinction between the Mass of the Catechumens and Mass of the Faithful identifies a real threshold, but it must not split the celebration into two unrelated services. Scripture culminates in the Gospel and Creed; the faith confessed becomes sacrificial offering; the Canon makes present the mystery proclaimed; Communion incorporates its recipients more deeply into what they have heard and offered. The Catechism therefore names two great parts that form one act: gathering and the liturgy of the Word, then the presentation of gifts, consecratory thanksgiving, and Communion (CCC 1345--1355; cf. \emph{Sacrosanctum Concilium} 48, 56).

The ordained celebrant and the faithful do not contribute interchangeable priesthoods. The ministerial priest acts sacramentally in the person of Christ the Head; the baptized exercise their common priesthood by uniting praise, self-offering, prayer, and worthy reception with Christ's oblation according to their condition (\emph{Lumen gentium} 10--11; \emph{Mediator Dei} 80--104). The recurring plurals of the Canon and the \latin{Orate, fratres} are therefore neither theatrical speech nor a democratic delegation of sacramental power. They disclose an ordered ecclesial action whose principal agent is Christ.

Finally, symbolic analysis remains sacramental realism. Bread, wine, altar, incense, washing, speech, silence, bow, kiss, fracture, and procession are not ciphers that replace what happens. They are bodily signs within an efficacious ecclesial rite. At the consecration, the signs do not merely awaken an idea of absence: by Christ's words and power, the substances of bread and wine become his Body and Blood while the sensible species remain (Trent, Session XIII, ch. 4; Session XXII, chs. 1--2; CCC 1373--1381).

Each numbered movement therefore begins with a fourfold source dossier. It first marks the stable Latin boundary or explains that a proper text varies; then returns biblical language to its literary and canonical setting; then gives a securely dated witness or a candid historical range; finally it reads Fathers, Aquinas, symbol, and adjacent units together. A dated decree, a probable manuscript range, and a retrospective attribution are not interchangeable evidence. This discipline also keeps biblical typology from becoming pseudo-history: Isaiah may illuminate a cleansing prayer without having prescribed its medieval placement.

\begin{massbox}{Scope limitation}
This is an exposition of an identified historical liturgical edition. It neither argues that every historical layer is immutable nor treats later reform as the interpretive purpose of the earlier book. Present permissions for celebrating with the 1962 Missal are a distinct, mutable canonical question and must be checked in the applicable universal and diocesan law. The final Leonine-prayer section concerns a law and devotion attached after Mass; it is outside the \latin{Ordo Missae}, outside the sacrificial action just analyzed, and outside the thirty-three numbered movements.
\end{massbox}
